\documentclass[10pt,a4paper]{amsart}
\usepackage[margin=19mm]{geometry}
\usepackage{amsmath,amssymb,mathtools}
\usepackage[scr=rsfs,cal=euler]{mathalfa}
\usepackage{xfrac}
\usepackage[unicode,hidelinks,bookmarks=false]{hyperref}
\newcommand{\ie}{i.e.\ }
\newcommand{\cf}{cf.\ }
\newcommand{\resp}{respectively\ }
\newcommand{\Subsection}{\subsection}
\providecommand{\nobreakdash}{\nobreak\hbox{-}}
\setlength{\emergencystretch}{2em}
\multlinegap0pt
\usepackage[shortlabels]{enumitem}
\usepackage{xcolor}
\usepackage{amssymb}
\newtheorem{lemma}{Lemma}
\newtheorem{claim}{Claim}
\title{On Balance, To What Degree is Burr's Conjecture True?}
\author{Shagnik Das, Bruce Reed, Jozef Skokan}
\date{Source: arXiv:2606.11410v1 (CC BY 4.0)}
\begin{document}
\maketitle

\noindent\emph{Source and licence.} On Balance, To What Degree is Burr's Conjecture True?. Version arXiv:2606.11410v1 (CC BY 4.0), distributed by arXiv under the Creative Commons Attribution 4.0 International licence, http://creativecommons.org/licenses/by/4.0/, as recorded in the arXiv licence metadata for this exact version and shown on the abstract page https://arxiv.org/abs/2606.11410v1. Full licence notice: \texttt{licenses/arxiv-2606.11410-CC-BY-4.0.txt}.\par\smallskip

\noindent\textit{Editorial scope.} Counted unit: the article's lemma commonlem (source0.tex lines 317--320). The article's complete original proof of it is delivered with it (source0.tex lines 321--442, one \texttt{proof} environment of 122 lines). No mathematics is written, repaired or re-derived: the statement and the proof are the article's own lines, in the article's own notation, and no retained line is substituted or re-flowed.  No further source-local context is needed: every cross-reference of every retained line resolves inside the two delivered spans.  Nothing was omitted from any retained span, so the package carries no omission record.  No retained line contains a citation, so the wrapper carries no bibliography and no bibliography entry.  No retained line contains a figure environment, an image-inclusion command or drawing code, so no image is delivered and the package declares no asset dependency.  Editorial formatting only, in addition: an AMS \texttt{amsart} preamble that restates the article's own declarations and macros used by the retained lines, namely the environments \texttt{lemma}, \texttt{claim} on separate counters, together with the standard packages those lines need.  The retained environments print in this document as Lemma 1, Claim 1 in order of appearance, whereas the article prints the same items with the numbering of its own full document.  The complete native source of the article as distributed in the arXiv e-print for this exact version is bundled unchanged as \texttt{sources/202548/source0.tex}.
\begin{lemma}
\label{commonlem}
    If there is no monochromatic copy of $T$ then, renaming colours if needed, we can partition $V$ into $C$ and $D$ so that $G_r[C,D]$, $G_b[C]$ and $G_b[D]$ all have minimum degree at least $\frac{n}{2}-2t_1$, while $\Delta(G_b[C,D]) \le t_1$. Furthermore, if $T$ has a center then $\Delta(G_r)<t_2$. 
\end{lemma}

\begin{proof}
We distinguish two cases.
   
\noindent{\bf Case 1:} $T$ has a center $c$.

Wlog $V-Z$ contains a vertex $v$ such that the degree of $v$  in $G_b$ is at least $\lceil n-1/2 \rceil \ge t_2-1$. Indeed either $v$ has degree at least $t_2$ or $n=2t_2-1$  and the degree of $v$ and every other vertex of $V-Z$ in $G_b$ is  $t_2-1$. Furthermore, in the second case, $\Delta(T)<t_2$.
Since $v$ is in $Z$ this  degree is at most $n/2+2t_1$.   If there 
are $2.01t_1t_2$ blue  edges between $N_b(v)$ and $V-N_b(v)-Z-v$ then $G_b[V-N_b(v)-v-Z,N_b(v)]$ has average degree exceeding $2t_1$ and contains a subgraph of minimum degree at least $t_1$ and average degree at least $2t_1$. We consider a maximal such subgraph of minimum degree $t_1$ and average degree $2t_1$ and note it  has at least $2t_1$ vertices on each side and so any vertex of $N_b(v)$ not contained in it is nonadjacent to more than $t_1$ vertices of its vertices outside $N_b(v)$. We embed $c$ in $v$ and $T_2-L_2-c$ in this subgraph with $I_1-c$ outside 
$N_b(v)$. Furthermore, if $v$ has degree $t_2-1$ in $G_b$ and there is a vertex $w$  in this subgraph which is not in $N_b(v)$ and which either has a blue edge to a vertex outside the subgraph or satisfies $|N_b(w) \cup N_b(v)|> t_2+t_1-3$
then we embed a vertex of $I_1-s$, all of whose children are leaves,  in $w$ as the first step of our embedding.  Since (i) every vertex of $G_b-Z$ has degree exceeding  $3t_2/4+t_1$, 
(ii) at least  $t_2/4$ vertices of $L_2$ are incident to $c$, and (iii) no vertices of $I_1$ are embedded in any of the at
least $t_2-1$ vertices of $N_b(v)$,  we can then embed $L_2$ by applying Hall's theorem to obtain a contradiction unless $|N_b(v)|=t_2-1$
and we could not find the desired $w$.We note that since, in this case,  all the vertices in the subgraph in $V-N_b(v)-Z$ have degree $t_2-1$ in $G_b$  this implies 
that for each such vertex $y$, $|N_b(v)-N_b(y)|<t_1$. Since our subgraph contains every vertex of $N_b(v)$ which is nonadjacent  in $G_b$ to 
fewer than $t_1$+1 vertices of the subgraph which are not in $N_b(v)$ this implies that the number of vertices 
in $N_b(v)$ which are not in the subgraph is less than the number of vertices outside $N_b(v)$ which are in the subgraph.
So every vertex of $G_b-Z$ sees a vertex of the subgraph. Now if there is a vertex $w'$ of  $V-Z-N_b(v)$  which is  not in the subgraph, then since we could not choose $w$,  $w'$ sees none of the subgraph outside $N_b(v)$ and hence sees a vertex of the subgraph in  $N_b(v)$. Furthermore, since it is not in the subgraph it sees at most $t_1-1$ such vertices and hence sees fewer than $t_2-t_1-2$   vertices in $N_b(v)$. 
Thus we can  embed a vertex of $I_1-s$ with only leaf children in $w'$, its parent in a vertex of the subgraph in $N_b(v)$ and proceed as before to 
to complete the embedding of $T$.  So, either there are at most $2.01t_1t_2$ blue edges between $N_b(v)$ and $V-N_b(v)-Z$ or every 
vertex of $V-N_b(v)-Z$ has blue edges to all but $t_1$ vertices of $N_b(v)$ and hence there are fewer than 
$2.01t_1t_2$ red edges between these two sets. Swapping colours if necessary and repeatedly deleting the  vertex  of maximum  degree in 
$G_b[V-N(v)-v-Z,N(v)-Z]$ until we have deleted $\sqrt{2.01t_2t_1}$ vertices we obtain a bipartite subgraph $G_r[C',D']$ of $G_r$  of  minimum degree $\frac{n}{2}-2\sqrt{2.01t_2t_1}-6t_1$ each of whose colour classes contains at least 
$\frac{n}{2}-\sqrt{2.01t_2t_1}-6t_1$ vertices.

If any vertex $v$ satisfies $|N_r(v)| \ge t_2,|N_r(v) \cap C'| \ge t_1,|N_r(v) \cap D'| \ge t_1$,  then we  choose a set $M$ of $t_2$  neighbours of $v$ in $G_r$ containing $t_1$ vertices of both $C'$ and $D'$.
Wlog we can insist that $M \cap C' \le \frac{t_2}{2}$ and $M \cap D' \ge \frac{t_2}{3}$. We embed $c$ in $v$ and  the roots of the nonsingleton components 
of $T-c$ in $D'$. We embed  the rest of $T_2-L_2$ in $G_r[C'-M,D']$ which implies  we embed none of $I_1$ in $M$. We can 
now embed  $L_2$ by applying Hall's Theorem. 

So any vertex $v$ satisfying  $|N_r(v)| \ge t_2$ satisfies either $|N_r(v) \cap C'| <t_1$ or $|N_r(v) \cap D'| < t_1$. 
In particular , $G_r[C']$ and $G_r[D']$ have maximum degree $2\sqrt{2.01t_2t_1}+6t_1$ and hence 
$G_b[C']$ and $G_b[D']$ have minimum degree $\frac{n}{2}-3\sqrt{2.01t_2t_1}-12t_1$. Furthermore, 
any vertex satisfying  $|N_r(v)| \ge t_2$ has at least $\frac{n}{2}-\sqrt{2.01t_2t_1}-7t_1$ blue edges to one of 
$C'$ or $D'$ and at least $t_2-2\sqrt{2.01t_2t_1}-13t_1$ red edges to the other. 


If any vertex $v$ satisfies $|N_b(v)| \ge \frac{n-1}{2},|N_b(v) \cap C'| \ge t_1,|N_b(v) \cap D'| \ge t_1$, then we can choose 
a set $M$ of $\min(t_2,\lceil \frac{n-1}{2}\rceil)$ neighbours of $v$ in $ G_b$ containing $t_1$ vertices of both $C'$ and $D'$.  Wlog $M$  contains at most $t_2/2$ vertices of  $C'$ and at least $t_2/3$ vertices of $D'$. We embed $c$ in $v$ and $T-L_2$ in $G[C']$ using $M$ only for the neighbours of $v$ which 
we embed first. We can then apply Hall's Theorem to embed $L_2$.

So any vertex $v$ satisfying  $|N_b(v)| \ge \frac{n-1}{2}$ satisfies either $|N_b(v) \cap C'| <t_1$ or $|N_b(v) \cap D'| < t_1$. 
Thus, any such   vertex has at least $\frac{n}{2}-\sqrt{2.01t_2t_1}-7t_1$ red  edges to one of 
$C'$ or $D'$ and at least $t_2 -2\sqrt{2.01t_2t_1}-13t_1$ blue  edges to the other. 

We let $C$ be those vertices which have more blue neighbours in $C'$ then $D'$ and $D=V-C$. 
We can now repeat the arguments of the first and third of the last four paragraphs with $C$ in place of $C'$ and $D$ in place of $D'$. 
We see that every vertex either (i) has degree at least $t_2$ in $G_r$  and at most $t_1$ red edges to the element of $(C,D)$,
containing it or (ii) has degree at least $\frac{n-1}{2}$ in $G_b$  and at most $t_1$ blue  edges to the element of $(C,D)$,
which does not contain it. We note that if  a vertex $v$  in $G-Z$ satisfies  (i) then we can embed $T-L_2$ with the center in $v$, $I_2-L_2$ in the side not containing $v$ and $I_1-v_1$ in nonneighbours of $v$ in the side containing it. We can then embed $L_2$. So, every vertex satisfies (ii) and not (i), and  letting $d$ be the maximum degree of $G_b(C,D)$ we have  $d \le t_1$ and that  
$\frac{n}{2}-d \le \min(G_b(C),G_b(D)) \le |C|,|D| < \frac{n}{2}+d $.
\vskip0.2cm


\noindent{\bf Case 2:} $T$ has no center.



\begin{claim}
 For some  $v\in V-Z$ fewer than  $\frac{n}{2}+8t_1$ vertices   have a red edge to  more than  
$3t_1$ blue-neighbours of $v$    
\end{claim}
\begin{proof}
    We assume the contrary and embed $T$ for a contradiction. We  first embed the at most $2t_1$ nodes of $T-L_2$ in $G_r-Z$, exploiting the fact that 
    this graph has minimum degree at least $\frac{n}{2}-6t_1$ to ensure  that 
whenever we embed  the second of the pair $(s_{2i-1},s_{2i})$, we embed it in a vertex which is 
adjacent to $3t_1$ blue-neighbours of 
the vertex in which the other element of the pair  is embedded. Since we have embedded at most $2t_1$ 
vertices, by  Hall's theorem  either we can now embed 
$L$ or there is a subset $W$ of $I_1$, such that the union of the neighbourhoods of 
the images of the elements of $W$ has size at most the minimum of $t_2+t_1$ and  $2t_1$ plus  the number of leaves 
incident to elements of $W$. By our choice of vertex images, this implies that for each $i$, 
$W$ contains at most one of $s_{2i-1}$ or $s_{2i}$. Since $T$ has no center this implies that 
the number of leaves adjacent to  the elements of $W$ is at most  $\frac{2t_2}{3}$. But the image of every 
element of $W$  has degree at least $\frac{n}{2}-t_1>\frac{2t_2}{3}+2t_1$ in $G_r$. so no such $W$ 
exists and we can embed $T$.
\end{proof}

So, we can and do choose  a vertex $v$ in $V-Z$ such that there are at most  
$\frac{n}{2}+8t_1$ vertices  which have a red edge to  more than $3t_1$ blue    neighbours of $v$.  
Hence there is set  $Y$ of at least $\frac{n}{2}-12t_1$ vertices of $V-Z$ each of which has red edges to at most $3t_1$
vertices of $V-N_r(v)$ and hence has red edges to at least 
$\frac{n}{2}-5t_1$ elements of $N_r(v)$. Any  vertex in $Y \cap N_r(v)$ has a red edge to   at most 
$3t_1$ vertices  of $Y-N_r(v)$. Any vertex of $Y-N_r(v)$ has a red edge 
to all but  at most $7t_1$ vertices of $N_r(v)$ and hence of $Y \cap N_r(v)$. So, either  $|Y \cap N_r(v)|\le 14t_1$
or $Y - N_r(v) \le 6t_1$.

If $|Y \cap N_r(v)| \le 14t_1$ vertices, then $| Y-N_r(v)| \ge \frac{n}{2}-26t_1$ and $G_r[Y-N_r(v),N_r(v)]$ 
is a bipartite graph, such that 
each vertex of $Y-N_r(v)$ is nonadjacent to at most $7t_1$ vertices of $N_r(v)$. Hence we can delete
the $\sqrt{7.1t_2t_1}$ vertices of $N_r(v)$ with lowest degree in this graph, to obtain a 
a bipartite subgraph  of $G_r$  which has minimum degree exceeding $n/2-\sqrt{7.1t_2t_1}-26t_1$.

If $Y - N_r(v)$ contains at most $6t_1$ vertices,$|Y \cap N_r(v)| \ge \frac{n}{2}-18t_1$ and   $ G_b [Y \cap N_r(v),V-N_r(v)]$ 
is a bipartite graph  such that 
each vertex of $Y \cap N_r(v)$ is nonadjacent to at most $3t_1$ vertices of $V-N(v)$. Hence we can delete
the $\sqrt{3.1t_2t_1}$ vertices of $V-N_r(v)$ with lowest degree in this graph, to obtain a 
a bipartite subgraph   of $G_r$ which has minimum degree exceeding $n/2-\sqrt{3.1t_2t_1}-18t_1$

So, wlog $G_r$ contains a bipartite subgraph $G_r[ C',D']$   which has minimum degree exceeding $n/2-\sqrt{7.1t_2t_1}-26t_1$.

If any vertex  has red edges to  at least $t_1$ vertices on one sode of this subgraph and at least $t_2/3$ on the other, then we can embed $T$ in $G_b$  as follows. We embed  $s$ in this vertex,  all the leaf 
children of $s$ on the  side in which the vertex  has  more red neighbours, and the remaining neighbours of $S$ so that  for $X \in \{C',D'\}$, the total number of vertices of $I_2$ 
in the components  containing neighbours of $S$  embedded  in $X$ is at least $\frac{t_2}{4}-t_1$.
We can then finish the embedding greedily in $G_r[C',D']$. So such $s$ exists and in particular 
 $G_b[C']$ and $G_b[D']$ have minimum degree exceeding $n/2-\sqrt{7.1t_2t_1}-27t_1$. %one place for 1/500

So,  if there is  a vertex $v$ in $V$ which has blue edges to at least $t_2/3$ vertices on one side of the 
subgraph and  at least $t_1$ vertices on the other   we can embed $T$ in $G_b$ as follows. 
We embed $s$ in this vertex, all the leaf 
children of $s$ on the  side in which the vertex  has  more blue  neighbours, and the remaining neighbours of $S$ so that  for $X \in \{C',D'\}$, the total number of vertices of $I_2$ 
in the components  containing neighbours of $S$  embedded  in $X$ is at least $\frac{t_2}{4}-t_1$.
We can then finish the embedding greedily in $G_b[C'] 
\cup G[D']$. So no such $s$ exists.

The last two paragraphs imply that every vertex has less than  $t_1$ blue edges to one of $C'$ 
and $D'$ and less than  $t_1$ red edges to the other. 
So, we can partition $V$ into $C$ and $D$ with $|C| \le |D|$ so that $G_r[C,D], G_b[C]$ and $G_b[D]$ all have minimum  degree at least  $n/2-\sqrt{3.1t_2t_1}-19t_1$.

We can now repeat the above arguments with $C$ and $D$ in the place of $C '$ and $D'$ to obtain that 
  $\Delta(G_b[C,D]), \Delta(G_r[B]),$ and $\Delta(G_r[D])$ are all less than $t_1$.  This implies $|D| \le \frac{n}{2}+t_1$, as otherwise we could embed $T$ greedily in $G_r[C,D]$ with $I_1 \in C$. So, $|C| \ge \frac{n}{2}-t_1$
and $\delta(G_b[C]), \delta(G_b[D]), \delta(G_r[C,D]) > \frac{n}{2}-2t_1$. 
\end{proof}

\end{document}
